\section{Introduction}
\label{sec:introduction}
Crohn's disease (CD) is a chronic inflammatory disorder of the gastrointestinal tract with substantial variation among patients \cite{torres2017crohn}. RNA sequencing (RNA-seq) measures RNA abundance in tissue samples and provides expression estimates for thousands of genes \cite{conesa2016survey}. Comparing intestinal expression profiles from patients with CD and comparison groups can reveal disease-associated patterns that may inform biomarker research \cite{montesi2026tentacles}. These associations do not, by themselves, establish causation or clinical usefulness.

Machine-learning classifiers can use expression measurements to distinguish CD from \emph{Not IBD} samples. Yet the number of genes typically exceeds the number of patients, and many genes are correlated. Feature selection therefore narrows the measurements to a \emph{gene signature}: a subset used for prediction or subsequent biological analysis. This creates two separate questions. \emph{Predictive generalisation} asks whether a model correctly classifies samples from a cohort excluded from training. \emph{Gene-signature reproducibility} asks whether fitting a comparable selection procedure separately to different patient cohorts yields similar gene lists. A model can generalise while selecting different genes from one dataset to another.

For example, two correlated genes may carry similar predictive information, so a selection method can retain either one depending on the patients analysed. Ein-Dor \emph{et al.} found markedly different yet similarly predictive signatures in breast cancer expression data \cite{eindor2005outcome}; related work showed that reliable agreement could require much larger discovery samples \cite{eindor2006thousands}. These studies motivate measuring gene agreement separately from prediction. Agreement is not proof of a causal biomarker, and its magnitude also depends on signature size.

Cohorts introduce further variation through differences in tissue composition, patient populations, sequencing protocols and class balance. Random train--test splits within one study can preserve such study-specific characteristics in both partitions. Cross-study evaluation instead tests models on cohorts not used for development \cite{bernau2014cross}. We use leave-one-cohort-out (LOCO) evaluation, training on three cohorts and testing on the fourth in turn.

Existing approaches differ in what they try to stabilise. LASSO selects genes through sparse regression coefficients \cite{tibshirani1996lasso}; Elastic Net also accommodates correlated predictors \cite{zou2005elastic}. Stability Selection favours genes repeatedly selected from subsamples \cite{meinshausen2010stability}, whereas StableMate targets predictors whose associations remain comparatively stable across environments \cite{deng2024stablemate}. TENTACLES combines feature-importance evidence from several learning algorithms into a consensus panel \cite{montesi2026tentacles}. Selection frequency, between-environment stability and agreement between algorithms are distinct properties; none alone guarantees replication across independent cohorts.

TENTACLES has already studied this Crohn's disease collection and compared biomarker-selection approaches \cite{montesi2026tentacles}. We instead examine how \emph{predictive transfer} relates to \emph{gene-set consistency} across LASSO, Elastic Net, Stability Selection, StableMate and TENTACLES, anchored against simple no-selection, univariate differential-expression and size-matched random-panel baselines. Our four analytical cohorts contain 689 samples from six BioProjects and share 9,983 genes. Alongside method-specific predictions, a fixed downstream classifier assesses all selected panels and baselines under identical prediction settings. We compare gene overlap against size-aware random-panel expectations, evaluate fixed-$K$ panel-size sensitivity, test sensitivity to source quality-control flags in D3, and review Reactome pathway agreement.
